% ============================================================
%  MTH636 – Game Theory | Lecture 31
%  Incomplete Information Games -- Extensive Form, Strategies,
%  Utility, Bayesian Nash Equilibrium
% ============================================================
\documentclass[12pt, a4paper]{article}

% ============================================================
%  MTH636 -- Game Theory
%  Notes Preamble — Color-Coded, TikZ-Enhanced
% ============================================================

% --- Core packages ---
\usepackage[T1]{fontenc}
\usepackage[utf8]{inputenc}
\usepackage{lmodern}
\usepackage[margin=2.5cm, top=2.8cm, bottom=2.8cm]{geometry}
\usepackage{microtype}

% --- Math ---
\usepackage{amsmath, amssymb, amsthm, mathtools}
\usepackage{bm}        % bold math
\usepackage{cancel}    % strike-through in math

% --- Color & Boxes ---
\usepackage[dvipsnames, svgnames, x11names, table]{xcolor}
\usepackage[most]{tcolorbox}
\tcbuselibrary{breakable, skins, theorems}

% --- TikZ & Diagrams ---
\usepackage{tikz}
\usepackage{tikz-qtree}
\usetikzlibrary{arrows.meta, positioning, shapes.geometric, fit, backgrounds,
                decorations.pathreplacing, calc, matrix, trees}
\usepackage{pgfplots}
\pgfplotsset{compat=1.18}

% --- Tables ---
\usepackage{booktabs}
\usepackage{multirow}
\usepackage{array}
\usepackage{tabularx}
\usepackage{makecell}

% --- Typography & Layout ---
\usepackage{enumitem}
\usepackage{parskip}
\usepackage{fancyhdr}
\usepackage{titlesec}
\usepackage{titletoc}
\usepackage{hyperref}
\usepackage{graphicx}
\usepackage{float}
\usepackage{caption}

% --- Color Palette ---
\definecolor{defblue}{HTML}{1A5276}      % Definition header
\definecolor{defbluebg}{HTML}{D6EAF8}    % Definition background
\definecolor{thmgreen}{HTML}{1E8449}     % Theorem header
\definecolor{thmbg}{HTML}{D5F5E3}        % Theorem background
\definecolor{exorange}{HTML}{B7770D}     % Example header
\definecolor{exbg}{HTML}{FEF9E7}         % Example background
\definecolor{warnred}{HTML}{922B21}      % Warning/Key Insight header
\definecolor{warnbg}{HTML}{FADBD8}       % Warning background
\definecolor{noticegray}{HTML}{2C3E50}   % Note header
\definecolor{noticebg}{HTML}{EBF5FB}     % Note background
\definecolor{taketawayyellow}{HTML}{7D6608} % Key Takeaway header
\definecolor{takeawaybg}{HTML}{FDFEFE}   % Takeaway background
\definecolor{accentpurple}{HTML}{6C3483} % Accent / proof
\definecolor{lightgray}{HTML}{F2F3F4}    % shading

% ============================================================
%  TCOLORBOX ENVIRONMENTS
% ============================================================

% -- Definition (Blue) --
\newtcolorbox{defbox}[2][]{
  enhanced, breakable,
  colback=defbluebg, colframe=defblue, coltitle=white,
  fonttitle=\bfseries\sffamily,
  title={Definition: #2},
  attach boxed title to top left={yshift=-2mm, xshift=4mm},
  boxed title style={colback=defblue, rounded corners},
  arc=4pt, boxrule=1.2pt,
  left=6pt, right=6pt, top=8pt, bottom=6pt,
  #1
}

% -- Theorem / Proposition (Green) --
\newtcolorbox{thmbox}[2][]{
  enhanced, breakable,
  colback=thmbg, colframe=thmgreen, coltitle=white,
  fonttitle=\bfseries\sffamily,
  title={Theorem: #2},
  attach boxed title to top left={yshift=-2mm, xshift=4mm},
  boxed title style={colback=thmgreen, rounded corners},
  arc=4pt, boxrule=1.2pt,
  left=6pt, right=6pt, top=8pt, bottom=6pt,
  #1
}

% -- Proposition (Green variant) --
\newtcolorbox{propbox}[2][]{
  enhanced, breakable,
  colback=thmbg, colframe=thmgreen!70!black, coltitle=white,
  fonttitle=\bfseries\sffamily,
  title={Proposition: #2},
  attach boxed title to top left={yshift=-2mm, xshift=4mm},
  boxed title style={colback=thmgreen!70!black, rounded corners},
  arc=4pt, boxrule=1.2pt,
  left=6pt, right=6pt, top=8pt, bottom=6pt,
  #1
}

% -- Example (Orange/Yellow) --
\newtcolorbox{exbox}[2][]{
  enhanced, breakable,
  colback=exbg, colframe=exorange, coltitle=white,
  fonttitle=\bfseries\sffamily,
  title={Example: #2},
  attach boxed title to top left={yshift=-2mm, xshift=4mm},
  boxed title style={colback=exorange, rounded corners},
  arc=4pt, boxrule=1.2pt,
  left=6pt, right=6pt, top=8pt, bottom=6pt,
  #1
}

% -- Key Insight / Warning (Red) --
\newtcolorbox{warnbox}[2][]{
  enhanced, breakable,
  colback=warnbg, colframe=warnred, coltitle=white,
  fonttitle=\bfseries\sffamily,
  title={Key Insight: #2},
  attach boxed title to top left={yshift=-2mm, xshift=4mm},
  boxed title style={colback=warnred, rounded corners},
  arc=4pt, boxrule=1.2pt,
  left=6pt, right=6pt, top=8pt, bottom=6pt,
  #1
}

% -- Note (Gray-Blue) --
\newtcolorbox{notebox}[1][]{
  enhanced, breakable,
  colback=noticebg, colframe=noticegray!60, coltitle=white,
  fonttitle=\bfseries\sffamily,
  title={Note},
  attach boxed title to top left={yshift=-2mm, xshift=4mm},
  boxed title style={colback=noticegray!60, rounded corners},
  arc=4pt, boxrule=0.8pt,
  left=6pt, right=6pt, top=8pt, bottom=6pt,
  #1
}

% -- Key Takeaway --
\newtcolorbox{takebox}[1][]{
  enhanced, breakable,
  colback=takeawaybg, colframe=taketawayyellow, coltitle=white,
  fonttitle=\bfseries\sffamily,
  title={Key Takeaway},
  attach boxed title to top left={yshift=-2mm, xshift=4mm},
  boxed title style={colback=taketawayyellow, rounded corners},
  arc=4pt, boxrule=1.2pt,
  left=6pt, right=6pt, top=8pt, bottom=6pt,
  #1
}

% -- Intuition --
\newtcolorbox{intubox}[1][]{
  enhanced, breakable,
  colback=accentpurple!8, colframe=accentpurple!60, coltitle=white,
  fonttitle=\bfseries\sffamily,
  title={Intuition},
  attach boxed title to top left={yshift=-2mm, xshift=4mm},
  boxed title style={colback=accentpurple!80, rounded corners},
  arc=4pt, boxrule=0.8pt,
  left=6pt, right=6pt, top=8pt, bottom=6pt,
  #1
}

% ============================================================
%  PAYOFF MATRIX HELPER
%  Usage: inside tabular, use \payoff{a}{b} for (a,b)
% ============================================================
\newcommand{\payoff}[2]{$(#1,\;#2)$}
\newcommand{\nash}[2]{\underline{$#1$},\underline{$#2$}}  % underline NE payoffs
\newcommand{\Nset}{\mathbb{N}}
\newcommand{\Rset}{\mathbb{R}}
\newcommand{\Sset}{\mathcal{S}}
\newcommand{\Ui}{u_i}

% ============================================================
%  SECTION FORMATTING
% ============================================================
\titleformat{\section}
  {\large\bfseries\sffamily\color{defblue}}
  {\thesection}{1em}{}[\titlerule]

\titleformat{\subsection}
  {\normalsize\bfseries\sffamily\color{thmgreen!80!black}}
  {\thesubsection}{1em}{}

\titleformat{\subsubsection}
  {\normalsize\itshape\sffamily\color{exorange!80!black}}
  {\thesubsubsection}{1em}{}

% ============================================================
%  HEADER / FOOTER
% ============================================================
\pagestyle{fancy}
\fancyhf{}
\fancyhead[L]{\sffamily\small\color{defblue} MTH636 -- Game Theory}
\fancyhead[R]{\sffamily\small\color{defblue} \leftmark}
\fancyfoot[C]{\sffamily\small\thepage}
\renewcommand{\headrulewidth}{0.4pt}

% ============================================================
%  HYPERREF SETUP
% ============================================================
\hypersetup{
  colorlinks=true,
  linkcolor=defblue,
  urlcolor=thmgreen,
  citecolor=accentpurple,
  bookmarksnumbered=true,
  pdftitle={MTH636 Game Theory Notes},
  pdfauthor={MTH636 Student Notes}
}

% ============================================================
%  MISC HELPERS
% ============================================================
\newcommand{\BR}{\text{BR}}        % Best Response
\newcommand{\NE}{\text{NE}}        % Nash Equilibrium
\newcommand{\SPNE}{\text{SPNE}}    % Subgame Perfect NE
\newcommand{\Dom}{\text{dom}}      % Dominance
\newcommand{\mixs}[1]{\sigma_{#1}} % Mixed strategy
\newcommand{\pures}[1]{s_{#1}}     % Pure strategy


% ============================================================
\begin{document}

% Title
\begin{center}
  {\LARGE\bfseries\sffamily\color{defblue} MTH636 -- Game Theory}\\[6pt]
  {\Large\sffamily Lecture 31: Strategies, Utility \& Bayesian Nash Equilibrium}\\[4pt]
\end{center}
\vspace{0.5cm}
\hrule height 1.5pt
\vspace{0.3cm}

% ============================================================
\section{Extensive Form Representation of Harsanyi Games}

Recall from Lecture 30 that every Harsanyi game of incomplete information can be written as an extensive form game with imperfect information and chance nodes.

\begin{thmbox}{Harsanyi's Representation Theorem (restated)}
Every Harsanyi game of incomplete information can be written as an \textbf{extensive form game} with imperfect information and chance nodes.
\end{thmbox}

\noindent For the extensive form representation:
\begin{itemize}
  \item The number of \textbf{information sets} for player $i$ is $|T_i|$ -- one information set per possible type of player $i$.
  \item The game will have \textbf{perfect recall} for all players. This means we may use \textbf{behavior strategies}.
\end{itemize}

% ============================================================
\section{Types of Strategies}

Since the extensive form game has perfect recall, three kinds of strategies are relevant:

\subsection{Pure Strategies}

\begin{defbox}{Pure Strategy in a Harsanyi Game}
A \textbf{pure strategy} for player $i$ is a function
\[
  s_i : T_i \;\longrightarrow\; \bigcup_{t_i \in T_i} A_i(t_i),
\]
where $s_i(t_i) \in A_i(t_i)$ for every $t_i \in T_i$.
\end{defbox}

\subsection{Mixed Strategies}

\begin{defbox}{Mixed Strategy in a Harsanyi Game}
A \textbf{mixed strategy} for player $i$ is a probability distribution over the pure strategies of player $i$.
\end{defbox}

\subsection{Behavior Strategies}

\begin{defbox}{Behavior Strategy in a Harsanyi Game}
A \textbf{behavior strategy} for player $i$ is a function
\[
  b_i : T_i \;\longrightarrow\; \bigcup_{t_i \in T_i} \Delta A_i(t_i),
\]
where $b_i(t_i) \in \Delta A_i(t_i)$ for every $t_i \in T_i$.

Here $\Delta A_i(t_i)$ denotes the set of probability distributions (mixed actions) over $A_i(t_i)$.
\end{defbox}

\begin{notebox}
  In the extensive form, a behavior strategy for player $i$ assigns, to each information set (i.e.\ each type $t_i \in T_i$), a probability distribution over the available actions at that information set:
  \[
    \tilde{s}_i : H_i \;\longrightarrow\; \bigcup_{h \in H_i} A(h), \quad \tilde{s}_i(h_i) \in A(h_i).
  \]
  This corresponds exactly to $b_i(t_i) \in \Delta A_i(t_i)$ in the type-space notation.
\end{notebox}

% ============================================================
\section{Utility in Harsanyi Games}

\subsection{Utility at a Type Profile (Ex Ante)}

Consider a behavior strategy profile $b = (b_1, \ldots, b_n)$.

For a type profile $t = (t_1, \ldots, t_n)$ and an action profile $a = (a_1, \ldots, a_n)$ with $a_i \in A_i(t_i)$ for all $i \in \{1,\ldots,n\}$, the \textbf{utility of player $i$ at type profile $t$} is:
\[
  V_i(t; b) = \sum_{\substack{a \,\in\, \bigtimes_{i \in N} A_i(t_i)}}
  b_1(t_1; a_1) \times b_2(t_2; a_2) \times \cdots \times b_n(t_n; a_n) \cdot u_i(t, a),
\]
where $b_j(t_j; a_j)$ denotes the probability assigned by player $j$'s behavior strategy to action $a_j$ when her type is $t_j$.

\begin{notebox}
This is the \textbf{ex-ante utility} at type profile $t$ -- the expected utility before the randomisation (i.e.\ before the mixed actions are realised), but after the type profile is fixed. It is called the utility at the \emph{ex-ante stage}, before the randomisation takes place.
\end{notebox}

\subsection{Total (Ex Ante) Utility}

The \textbf{total (ex-ante) utility} of player $i$ under behavior strategy profile $b$ is:
\[
  V_i(b) = \sum_{t \in T} p(t)\, V_i(t; b).
\]

\subsection{Interim Utility}

After player $i$ knows her own type $t_i$ (but before learning the types of others), she can compute her \textbf{interim (conditional) utility}:
\[
  U_i(b \mid t_i) = \sum_{t_{-i} \in T_{-i}} p(t_{-i} \mid t_i)\, U_i\!\left((t_i, t_{-i}),\, b\right).
\]
This is called player $i$'s \textbf{interim utility} -- she knows her own type $t_i$ and takes expectation over the opponents' types using the conditional belief $p(t_{-i} \mid t_i)$.

\subsection{Relationship Between Ex-Ante and Interim Utility}

The following chain of equalities shows that the ex-ante utility can be decomposed through the interim utility:

\begin{align*}
  V_i(b)
  &= \sum_{t \in T} p(t)\, V_i(t; b) \\[4pt]
  &= \sum_{t \in T} p(t_{-i} \mid t_i)\, p(t_i)\, V_i(t; b) \\[4pt]
  &= \sum_{t_i \in T_i} \sum_{t_{-i} \in T_{-i}} p(t_i) \Bigl( p(t_{-i} \mid t_i)\, V_i(t; b) \Bigr) \\[4pt]
  &= \sum_{t_i \in T_i} p(t_i) \sum_{t_{-i} \in T_{-i}} p(t_{-i} \mid t_i)\, V_i(t; b) \\[4pt]
  &= \sum_{t_i \in T_i} p(t_i)\, U\!\left(b \mid t_i\right).
\end{align*}

\begin{takebox}
  The ex-ante utility is a weighted sum of interim utilities:
  \[
    V_i(b) = \sum_{t_i \in T_i} p(t_i)\, U_i(b \mid t_i).
  \]
  This is analogous to the law of total expectation.
\end{takebox}

% ============================================================
\section{Nash Equilibrium and Bayesian Nash Equilibrium}

\subsection{Nash Equilibrium (NE)}

\begin{defbox}{Nash Equilibrium in a Harsanyi Game}
A strategy profile $\sigma^* = (\sigma_i)_{i \in N}$ is a \textbf{Nash Equilibrium (NE)} if
\[
  V_i(\sigma^*) \geq V_i(\sigma_i,\, \sigma_{-i}^*) \quad \forall\, i \in N \text{ and all } \sigma_i \in B_i,
\]
where $B_i$ denotes the set of all behavior strategies for player $i$.
\end{defbox}

\subsection{Bayesian Nash Equilibrium (BNE)}

The Nash Equilibrium above is the \emph{ex-ante} equilibrium concept. A more refined concept -- the \textbf{Bayesian Nash Equilibrium} -- requires optimality \emph{at the interim stage}, i.e., for every type the player might be.

\begin{defbox}{Bayesian Nash Equilibrium (BNE)}
A strategy profile $\sigma^*$ is a \textbf{Bayesian Nash Equilibrium (BNE)} if
\[
  U_i(\sigma^* \mid t_i) \geq U_i\!\left((a_i, \sigma_{-i}^*) \mid t_i\right)
  \quad \forall\, i \in N,\ \forall\, t_i \in T_i,\ \forall\, a_i \in A_i(t_i).
\]
That is, for every player $i$, for every type $t_i$, and for every possible deviation $a_i$, the equilibrium strategy $\sigma_i^*$ is at least as good as any deviation -- when evaluated using the interim (conditional) utility given type $t_i$.
\end{defbox}

\begin{notebox}
In the case where $T_i$ has multiple elements, say $t_1, t_{n}, t_2$ (as in an example with three types), the total utility decomposes as:
\[
  U_i(\sigma) = p(t_1)\, U_i(\sigma \mid t_1) + p(t_n)\, U_i(\sigma \mid t_n) + p(t_2)\, U_i(\sigma \mid t_2).
\]
\end{notebox}

% ============================================================
\section{Key Results}

\begin{propbox}{Every BNE is also an NE}
Every Bayesian Nash Equilibrium is also a Nash Equilibrium.
\end{propbox}
\begin{proof}
See Solan's book.
\end{proof}

\begin{thmbox}{Existence of NE in Behavior Strategies}
For an incomplete information game (in the Harsanyi sense), a Nash Equilibrium exists in \textbf{behavior strategies}.
\end{thmbox}
\begin{proof}
The corresponding extensive form game has a Nash Equilibrium in mixed strategies. Since the game has perfect recall, every mixed strategy is equivalent to a behavior strategy. Therefore a behavior strategy Nash Equilibrium exists.
\end{proof}

\begin{intubox}
  The key steps in the existence proof are:
  \begin{enumerate}
    \item The Harsanyi game $\Rightarrow$ extensive form with chance nodes and information sets.
    \item The extensive form game has a Nash Equilibrium in mixed strategies (by Nash's theorem, since the game is finite).
    \item Perfect recall $\Rightarrow$ every mixed strategy is equivalent to a behavior strategy (Kuhn's theorem).
    \item Therefore a \emph{behavior strategy} NE exists.
  \end{enumerate}
\end{intubox}

\end{document}
